
Comme @@M0@@, on a @@M1@@, et par suite :
@@M2@@
Posons :
@@M3@@
Comme @@M4@@, on a :
@@M5@@
On en déduit, d'après V (6.2.1) :
@@M6@@
Compte tenu de 3.5, on obtient finalement :
@@M7@@

D'après la définition de @@M8@@, on a donc :
@@M9@@

\textbf{5.4.} En utilisant la base de @@M10@@ sur @@M11@@ formée des éléments @@M12@@, avec @@M13@@, on voit par additivité qu'il suffit de prouver le théorème pour des éléments de la forme @@M14@@. On a alors :
@@M15@@
D'autre part :
@@M16@@

On voit donc qu'il suffit de montrer le théorème pour les @@M17@@, avec @@M18@@. Or d'après (VI (5.2.3)), on a
@@M19@@
Par suite :
@@M20@@

